\documentclass[11pt,a4paper]{article}
\usepackage[utf8]{inputenc}
\usepackage[french]{babel}
\usepackage[T1]{fontenc}
\usepackage{geometry}
\usepackage{enumitem}
\usepackage{titlesec}
\usepackage{hyperref}
\usepackage{xcolor}
\usepackage{longtable}

\geometry{top=2cm, bottom=2cm, left=2.5cm, right=2.5cm}

\titleformat{\section}
{\normalfont\Large\bfseries\color{green!70!black}}{\thesection}{1em}{}

\titleformat{\subsection}
{\normalfont\large\bfseries\color{green!50!black}}{\thesubsection}{1em}{}

\title{\textbf{Guide Développeur 2\\E-Salle ENSAA}}
\author{Gestion Avancée}
\date{\today}

\begin{document}

\maketitle
\tableofcontents
\newpage

\section{Vue d'Ensemble - Développeur 2}

\subsection{Votre Mission}

Vous êtes responsable de la \textbf{gestion avancée} du système E-Salle ENSAA avec des modules complexes.

\textbf{Vos 4 modules :}
\begin{enumerate}
    \item \textbf{Administration} (5 jours) - Semaine 1
    \item \textbf{Matière} (8 jours) - Semaines 2-3
    \item \textbf{Emploi du Temps} (10 jours) - Semaines 4-5
    \item \textbf{Réservation} (6 jours) - Semaine 6
\end{enumerate}

\textbf{Total :} 29 jours de développement pur

\subsection{Particularités de Votre Position}

\begin{itemize}
    \item ⚠️ \textbf{1 dépendance :} Matière attend Filière de Dev1 (fin S3)
    \item ✅ \textbf{Modules complets} : Backend + Frontend pour chaque module
    \item ⚠️ \textbf{Modules complexes} : Logique métier avancée, algorithmes, UI interactive
    \item ✅ \textbf{Aide de Dev1} : Dev1 vous aide en S5-S6 sur Emploi et Réservation
\end{itemize}

\subsection{Gestion de l'Attente (Semaine 3)}

\textbf{Stratégie :}
\begin{itemize}
    \item S1 : Faites Administration (ne dépend que de Users)
    \item S2 : Continuez Administration + préparez Matière (UI, structure)
    \item S3 (début) : Attendez que Dev1 termine Filière (3 jours)
    \item S3 (fin) : Commencez Matière dès que Filière est mergée (3 jours)
\end{itemize}

\newpage

\section{Vos Modules en Détail}

\subsection{Module 1 : Administration (Semaine 1 - 5 jours)}

\textbf{Objectif :} Système d'approbation des inscriptions et gestion des utilisateurs.

\textbf{Fonctionnalités :}
\begin{enumerate}
    \item \textbf{Dashboard admin} : Liste des inscriptions EN\_ATTENTE
    \item \textbf{Approuver inscription} :
    \begin{itemize}
        \item Si PROFESSEUR/COORDINATEUR : Sélectionner filière → Statut = ACTIF
        \item Si MEMBRE\_CLUB : Saisir nom du club → Statut = ACTIF
    \end{itemize}
    \item \textbf{Refuser inscription} : Statut = REFUSE
    \item \textbf{Notifications} : Email + WhatsApp au user (approbation/refus)
    \item \textbf{Gestion users} : Liste, modifier, supprimer
\end{enumerate}

\textbf{Règles métier :}
\begin{itemize}
    \item Seul l'admin peut approuver/refuser
    \item Notification immédiate au user lors de l'approbation/refus
    \item Pour PROFESSEUR/COORDINATEUR : filiere\_id obligatoire
    \item Pour MEMBRE\_CLUB : nom\_club obligatoire (texte libre)
\end{itemize}

\textbf{Fichiers à créer :}
\begin{itemize}
    \item \texttt{AdminServlet.java} (approbation, refus, gestion users)
    \item \texttt{admin/dashboard.jsp} (liste inscriptions EN\_ATTENTE)
    \item \texttt{admin/approve-form.jsp} (formulaire approbation avec sélection filière/club)
    \item \texttt{admin/users-list.jsp} (gestion users)
\end{itemize}

\subsection{Module 2 : Matière (Semaines 2-3 - 8 jours)}

\textbf{⚠️ DÉPENDANCE :} Attend Filière de Dev1 (fin S3)

\textbf{Objectif :} Gestion des matières avec charges horaires et professeurs.

\textbf{Entités :}
\begin{itemize}
    \item \texttt{Matiere} : id, nom, filiere\_id, professeur\_id, heures\_cours, heures\_td, heures\_tp, coefficient
\end{itemize}

\textbf{Fonctionnalités :}
\begin{enumerate}
    \item \textbf{CRUD complet} : Créer, Lire, Modifier, Supprimer
    \item \textbf{Gestion charges horaires} : Cours, TD, TP (séparés)
    \item \textbf{Assignation professeur} : Sélection d'un user avec role PROFESSEUR ou COORDINATEUR
    \item \textbf{Assignation filière} : Sélection d'une filière existante
\end{enumerate}

\textbf{Règles métier :}
\begin{itemize}
    \item Nom unique par filière
    \item Au moins 1 type d'heures (cours, TD ou TP) > 0
    \item Professeur doit avoir role = PROFESSEUR ou COORDINATEUR
    \item Filière doit exister (FK vers Filiere de Dev1)
\end{itemize}

\textbf{Fichiers à créer :}
\begin{itemize}
    \item \texttt{Matiere.java} (domain)
    \item \texttt{MatiereRepository.java + MatiereRepositoryImpl.java}
    \item \texttt{MatiereService.java + MatiereServiceImpl.java}
    \item \texttt{MatiereServlet.java}
    \item \texttt{list.jsp}, \texttt{form.jsp}, \texttt{view.jsp} (dans \texttt{views/matiere/})
\end{itemize}

\subsection{Module 3 : Emploi du Temps (Semaines 4-5 - 10 jours)}

\textbf{Objectif :} Interface calendrier interactive pour créer les emplois du temps.

\textbf{Entités :}
\begin{itemize}
    \item \texttt{EmploiDuTemps} : id, filiere\_id, annee, semaine\_debut, semaine\_fin
    \item \texttt{Seance} : id, emploi\_id, matiere\_id, salle\_id, professeur\_id, jour, heure\_debut, heure\_fin, groupe (1/2), statut (ACTIVE/LIBEREE)
\end{itemize}

\textbf{Fonctionnalités :}
\begin{enumerate}
    \item \textbf{Interface calendrier} :
    \begin{itemize}
        \item Onglets par année (DLA1, DLA2, DLA3 pour Cycle Ingénieur)
        \item Grille : Jours (Lun-Sam) × Créneaux horaires (8h-18h)
        \item Bouton "+" dans chaque créneau vide
    \end{itemize}
    \item \textbf{Modal d'ajout séance} :
    \begin{itemize}
        \item Sélection matière (filtrée par filière)
        \item Sélection salle (filtrée par type : COURS pour cours, TP pour TP, TD pour TD)
        \item Sélection professeur (filtrée par matière)
        \item Sélection groupe (1 ou 2) si TP/TD
    \end{itemize}
    \item \textbf{Génération 7 séances} : Répéter la séance pour 7 semaines consécutives
    \item \textbf{Libération exceptionnelle} : Professeur peut libérer 1 séance spécifique
    \item \textbf{Libération définitive} : Professeur peut libérer toutes les séances d'une matière
\end{enumerate}

\textbf{Règles métier :}
\begin{itemize}
    \item Pas de chevauchement : 1 salle = 1 séance par créneau
    \item Pas de chevauchement : 1 professeur = 1 séance par créneau
    \item Salles filtrées par type (TP uniquement pour TP)
    \item Groupe obligatoire si TP/TD
    \item Notification au coordinateur et admin lors de libération
\end{itemize}

\textbf{Fichiers à créer :}
\begin{itemize}
    \item \texttt{EmploiDuTemps.java}, \texttt{Seance.java} (domain)
    \item \texttt{EmploiDuTempsRepository.java + Impl}, \texttt{SeanceRepository.java + Impl}
    \item \texttt{EmploiDuTempsService.java + Impl}, \texttt{SeanceService.java + Impl}
    \item \texttt{EmploiDuTempsServlet.java}, \texttt{SeanceServlet.java}
    \item \texttt{calendar.jsp} (interface calendrier avec onglets)
    \item \texttt{modal-add-seance.jsp} (modal pour ajouter séance)
    \item JavaScript pour l'interactivité (bouton +, modal, AJAX)
\end{itemize}

\textbf{⚠️ Module le plus complexe !} Dev1 vous aide en S5 !

\subsection{Module 4 : Réservation (Semaine 6 - 6 jours)}

\textbf{Objectif :} Système de réservation avec priorités, FIFO, et vérifications.

\textbf{Entités :}
\begin{itemize}
    \item \texttt{Reservation} : id, user\_id, salle\_id, date, heure\_debut, heure\_fin, motif, priorite, statut (EN\_ATTENTE/APPROUVEE/REFUSEE), date\_demande
\end{itemize}

\textbf{Fonctionnalités :}
\begin{enumerate}
    \item \textbf{Demande de réservation} : User sélectionne salle, date, créneau, motif
    \item \textbf{Vérifications automatiques} :
    \begin{itemize}
        \item Salle disponible (pas dans Emploi du Temps)
        \item Salle disponible (pas dans autres Réservations)
        \item Salle TP bloquée pour MEMBRE\_CLUB
    \end{itemize}
    \item \textbf{Système de priorités} :
    \begin{itemize}
        \item PROFESSEUR/COORDINATEUR : priorité = 100
        \item MEMBRE\_CLUB : priorité = 50
        \item Si conflit : FIFO (First In First Out) dans chaque catégorie
    \end{itemize}
    \item \textbf{Alternatives} : Si salle indisponible, proposer salles similaires (même capacité, même type)
    \item \textbf{Notifications} : Email + WhatsApp (confirmation, refus, alternatives)
\end{enumerate}

\textbf{Règles métier :}
\begin{itemize}
    \item Double vérification : Emploi du Temps + Réservations existantes
    \item Salles TP bloquées pour MEMBRE\_CLUB
    \item Priorité PROFESSEUR > MEMBRE\_CLUB
    \item FIFO dans chaque catégorie de priorité
    \item Notification immédiate après décision
\end{itemize}

\textbf{Fichiers à créer :}
\begin{itemize}
    \item \texttt{Reservation.java} (domain)
    \item \texttt{ReservationRepository.java + ReservationRepositoryImpl.java}
    \item \texttt{ReservationService.java + ReservationServiceImpl.java} (logique priorités, FIFO, vérifications)
    \item \texttt{ReservationServlet.java}
    \item \texttt{list.jsp}, \texttt{form.jsp}, \texttt{view.jsp} (dans \texttt{views/reservation/})
\end{itemize}

\textbf{⚠️ Module complexe !} Dev1 vous aide si besoin en S6 !

\newpage

\section{Stratégie Git/GitHub pour Dev2}

\subsection{Configuration Initiale (Semaine 0)}

\textbf{Branches principales :}
\begin{itemize}
    \item \texttt{main} : Production (ne JAMAIS push directement)
    \item \texttt{develop} : Intégration (ne JAMAIS push directement)
    \item \texttt{dev2} : VOTRE branche de travail
\end{itemize}

\textbf{Setup initial :}
\begin{verbatim}
git clone <repo-url>
git checkout -b dev2
git push -u origin dev2
\end{verbatim}

\subsection{Workflow Quotidien}

\textbf{Chaque matin (avant de commencer) :}
\begin{verbatim}
git checkout dev2
git pull origin develop  # Récupérer les changements de Dev1
git pull origin dev2     # Récupérer vos propres changements
\end{verbatim}

\textbf{Pendant le développement :}
\begin{verbatim}
# Commits fréquents (2-3 fois par jour)
git add .
git commit -m "feat(admin): add user approval workflow"
git push origin dev2
\end{verbatim}

\textbf{Conventions de commit :}
\begin{itemize}
    \item \texttt{feat(module): description} - Nouvelle fonctionnalité
    \item \texttt{fix(module): description} - Correction de bug
    \item \texttt{refactor(module): description} - Refactoring
    \item \texttt{docs(module): description} - Documentation
\end{itemize}

\textbf{Exemples pour vos modules :}
\begin{verbatim}
feat(admin): add inscription approval
feat(matiere): add CRUD operations
feat(emploi): add calendar interface
feat(reservation): add priority system
fix(emploi): fix seance overlap validation
\end{verbatim}

\subsection{Workflow Hebdomadaire}

\textbf{Chaque vendredi (fin de semaine) :}

\textbf{1. Préparer votre code :}
\begin{verbatim}
git checkout dev2
git pull origin develop  # S'assurer d'avoir les derniers changements
# Résoudre conflits si nécessaire
git add .
git commit -m "merge: integrate develop changes"
git push origin dev2
\end{verbatim}

\textbf{2. Créer Pull Request (PR) :}
\begin{itemize}
    \item Aller sur GitHub
    \item Créer PR : \texttt{dev2 → develop}
    \item Titre : \texttt{[S1] Module Administration - Dev2}
    \item Description : Lister les fonctionnalités ajoutées
    \item Demander review à Dev1
\end{itemize}

\textbf{3. Après approbation :}
\begin{itemize}
    \item Merger la PR sur \texttt{develop}
    \item Dev1 pourra utiliser votre code (si besoin)
\end{itemize}

\subsection{Points de Synchronisation Critiques}

\textbf{Fin Semaine 1 :} Merger Administration → Module terminé

\textbf{Fin Semaine 3 :} Attendre que Dev1 merge Filière ⚠️ \textbf{CRITIQUE !}

\textbf{Après merge Filière :} Commencer Matière immédiatement

\textbf{Fin Semaine 4 :} Merger Matière → Vous pouvez continuer Emploi du Temps

\textbf{Fin Semaine 5 :} Merger Emploi du Temps → Vous pouvez commencer Réservation

\textbf{Fin Semaine 6 :} Merger Réservation → Tous vos modules sont terminés !

\subsection{Gestion des Conflits}

\textbf{Si conflit lors du pull :}
\begin{verbatim}
git pull origin develop
# Git affiche les conflits
# Ouvrir les fichiers en conflit
# Chercher <<<<<<< HEAD
# Choisir la bonne version ou combiner
git add <fichiers-résolus>
git commit -m "merge: resolve conflicts with develop"
git push origin dev2
\end{verbatim}

\textbf{Conseil :} Communiquez avec Dev1 quotidiennement pour éviter les conflits !

\subsection{Collaboration avec Dev1 (S5-S6)}

\textbf{Si Dev1 vous aide sur votre branche :}
\begin{verbatim}
# Dev1 travaille sur dev2
# Vous devez pull régulièrement
git checkout dev2
git pull origin dev2  # Récupérer les changements de Dev1
\end{verbatim}

\textbf{Communication :} Coordonnez-vous pour ne pas travailler sur les mêmes fichiers en même temps !

\newpage

\section{Checklist de Succès}

\subsection{Pour Chaque Module}

\textbf{Avant de merger :}
\begin{itemize}
    \item[$\square$] Entité créée avec annotations Hibernate
    \item[$\square$] Repository implémenté (CRUD complet)
    \item[$\square$] Service implémenté avec logique métier complexe
    \item[$\square$] Servlet créé avec toutes les routes
    \item[$\square$] Pages JSP créées (list, form, view)
    \item[$\square$] Validation des données (Bean Validation)
    \item[$\square$] Gestion des erreurs (try-catch)
    \item[$\square$] Notifications implémentées (Email + WhatsApp)
    \item[$\square$] Tests manuels effectués
    \item[$\square$] Code compilé sans erreur
    \item[$\square$] Commit + Push sur dev2
    \item[$\square$] PR créée vers develop
\end{itemize}

\subsection{Points de Synchronisation}

\textbf{Fin Semaine 1 :}
\begin{itemize}
    \item[$\square$] Module Administration terminé et mergé
    \item[$\square$] Workflow approbation/refus fonctionne
\end{itemize}

\textbf{Fin Semaine 3 :}
\begin{itemize}
    \item[$\square$] Filière de Dev1 mergée ⚠️ \textbf{CRITIQUE !}
    \item[$\square$] Module Matière commencé
\end{itemize}

\textbf{Fin Semaine 4 :}
\begin{itemize}
    \item[$\square$] Module Matière terminé et mergé
    \item[$\square$] Emploi du Temps en cours
\end{itemize}

\textbf{Fin Semaine 5 :}
\begin{itemize}
    \item[$\square$] Module Emploi du Temps terminé et mergé
    \item[$\square$] Interface calendrier fonctionne
\end{itemize}

\textbf{Fin Semaine 6 :}
\begin{itemize}
    \item[$\square$] Module Réservation terminé et mergé
    \item[$\square$] Système de priorités fonctionne
    \item[$\square$] Tous vos modules sont terminés !
\end{itemize}

\subsection{Communication}

\textbf{Quotidien :}
\begin{itemize}
    \item[$\square$] Standup 10 min avec Dev1 (matin)
    \item[$\square$] Signaler tout blocage immédiatement
\end{itemize}

\textbf{Hebdomadaire :}
\begin{itemize}
    \item[$\square$] Revue de code croisée (vendredi)
    \item[$\square$] PR mergée avant le weekend
\end{itemize}

\section{Conseils Pratiques}

\begin{enumerate}
    \item \textbf{Gérez l'attente (S3) :} Profitez-en pour peaufiner Administration et préparer Matière (UI, structure)
    \item \textbf{Modules complexes :} Décomposez en petites tâches (1 jour max par tâche)
    \item \textbf{Emploi du Temps :} Commencez par l'UI statique, puis ajoutez l'interactivité
    \item \textbf{Réservation :} Testez les algorithmes de priorité et FIFO avec des cas réels
    \item \textbf{Demandez de l'aide :} Dev1 est là pour vous aider en S5-S6 !
    \item \textbf{Testez souvent :} Lancez l'application après chaque fonctionnalité
    \item \textbf{Commitez fréquemment :} 2-3 commits par jour minimum
    \item \textbf{Communiquez :} Signalez tout problème immédiatement
\end{enumerate}

\textbf{Bon développement ! 🚀}

\end{document}

